% !TEX root = ../main.tex

\documentclass[../main.tex]{subfiles}

\graphicspath{{\subfix{../images/}}}

\begin{document}

\subsection{Протокол MAVLink и сообщение RAW\_IMU}

MAVLink (\textit{Micro Air Vehicle Link})~--- лёгкий бинарный протокол
сериализации для беспилотных систем, разработанный Лоренцем Майером
в 2009~году~\cite{mavlink_docs}. Протокол широко используется в экосистемах
PX4 и ArduPilot~\cite{px4_docs}. Для получения сырых показаний датчиков
применяется сообщение \texttt{RAW\_IMU} (идентификатор~27 в диалекте
\texttt{common}), содержащее девять целочисленных полей в единицах внутренней
шкалы датчика (LSB~--- \textit{Least Significant Bit}):


\begin{table}[H]
\centering
\caption{Поля сообщения MAVLink \texttt{RAW\_IMU}}
\label{tab:rawimu}
\begin{tabularx}{\textwidth}{|>{\centering\arraybackslash}p{4.5cm}|
                               >{\centering\arraybackslash}p{3cm}|
                               >{\centering\arraybackslash}X|}
\hline
\textbf{Поле} & \textbf{Датчик} & \textbf{Описание} \\
\hline
\texttt{xacc, yacc, zacc}   & Акселерометр & Ускорение по осям X, Y, Z (LSB) \\
\hline
\texttt{xgyro, ygyro, zgyro}& Гироскоп     & Угловая скорость по осям (LSB) \\
\hline
\texttt{xmag, ymag, zmag}   & Магнитометр  & Магнитное поле по осям (LSB) \\
\hline
\end{tabularx}
\end{table}


\subsection{Модуль первичного считывания данных}	

На первом этапе работы был написан автономный модуль 
\texttt{raw\_data.cpp}, позволивший проверить
корректность соединения с Pixhawk и собрать сырые данные для анализа и
подбора калибровочных коэффициентов. Взаимодействие реализовано через
последовательный порт \texttt{/dev/ttyACM0} на скорости 115\,200~бод.
Инициализация порта выполнена через POSIX-интерфейс
\texttt{termios}~\cite{posix_termios}: 8~бит данных, без паритета, один
стоп-бит, без аппаратного управления потоком, таймаут~1~с.
 
\begin{lstlisting}[caption={Настройка последовательного порта через termios}]
int setup_serial_port(const char* port, int baudrate) {
    int fd = open(port, O_RDWR | O_NOCTTY);
    struct termios tty;
    memset(&tty, 0, sizeof(tty));
    tcgetattr(fd, &tty);
    cfsetospeed(&tty, baudrate);
    cfsetispeed(&tty, baudrate);
    tty.c_cflag |= (CLOCAL | CREAD);
    tty.c_cflag &= ~CSIZE;
    tty.c_cflag |= CS8;
    tty.c_cflag &= ~(PARENB | CSTOPB | CRTSCTS);
    tty.c_lflag &= ~(ICANON | ECHO | ECHOE | ISIG);
    tty.c_iflag &= ~(IXON | IXOFF | IXANY);
    tty.c_oflag &= ~OPOST;
    tty.c_cc[VMIN]  = 0;
    tty.c_cc[VTIME] = 10;
    tcsetattr(fd, TCSANOW, &tty);
    return fd;
}
\end{lstlisting}
 
MAVLink-парсер реализован в виде конечного автомата.

\texttt{mavlink\_parse\_char()} возвращает ненулевое значение только
при получении полного валидного пакета с проверкой CRC-16~\cite{mavlink_proto}.
Каждый принятый пакет \texttt{RAW\_IMU} записывается в CSV-файл с меткой
времени с точностью до микросекунды~(\texttt{gettimeofday()}); принудительный
сброс буфера выполняется каждые 50~пакетов.


\subsection{Результаты первичного считывания}
 
Сообщения \texttt{RAW\_IMU} стабильно поступали с частотой порядка
100~Гц, что соответствует штатной конфигурации IMU на
Pixhawk~\cite{px4_docs}. Сырые данные акселерометра в покое показали
устойчивое значение по оси~Z (проекция силы тяжести), а гироскоп~---
дрейф нуля, подлежащий компенсации на следующем этапе.


\end{document}

